\documentclass[10pt, letterpaper]{article}
% guidance from wenbin: 简历一般是两页，写的。 课程pj写成research exp, selected project, 写两页有更多的research, 多样性，
% 两个邮箱都放，shanghaitech先
% online preference，先写Objective, 一句话概括要干什么，黑体，包括毕设，包括course project， 再基于目标把干了什么事情按照逻辑排出来，数字化成果，黑体。除了paper核心，实验，还要说写了paper, 结果是什么，可以写很多行。第二篇写出来是IEEE trans, in submittions to iclr, 多少分， submission to xxx，L4DC写全称，
% gre对点看，建议考320以上就行了， 托福gre不用放cv
% 马上当个ta可以写上去，不是那么重要
% 还有俩项目都写上去
% 需要建个个人主页吗？ better，弄不好就算了，paper, github, 主页， poster
% sirui learning to dalegate， 海报放个链接在主页上 ***
% 细节： research interest给几个词就行了+一句话, 这些话更像文书里的
% Packages:
\usepackage[
    ignoreheadfoot, % set margins without considering header and footer
    top=1.5 cm, % seperation between body and page edge from the top
    bottom=1.5 cm, % seperation between body and page edge from the bottom
    left=1.7 cm, % seperation between body and page edge from the left
    right=1.7 cm, % seperation between body and page edge from the right
    footskip=1.0 cm, % seperation between body and footer
    % showframe % for debugging 
]{geometry} % for adjusting page geometry
\usepackage{titlesec} % for customizing section titles
\usepackage{tabularx} % for making tables with fixed width columns
\usepackage{array} % tabularx requires this
\usepackage[dvipsnames]{xcolor} % for coloring text
\definecolor{primaryColor}{RGB}{0, 0, 0} % define primary color
\usepackage{enumitem} % for customizing lists
\usepackage{fontawesome5} % for using icons
\usepackage{amsmath} % for math
\usepackage[
    pdftitle={Tianyue Zhou's CV},
    pdfauthor={Tianyue Zhou},
    pdfcreator={LaTeX with RenderCV},
    colorlinks=true,
    urlcolor=primaryColor
]{hyperref} % for links, metadata and bookmarks
\usepackage[pscoord]{eso-pic} % for floating text on the page
\usepackage{calc} % for calculating lengths
\usepackage{bookmark} % for bookmarks
\usepackage{lastpage} % for getting the total number of pages
\usepackage{changepage} % for one column entries (adjustwidth environment)
\usepackage{paracol} % for two and three column entries
\usepackage{ifthen} % for conditional statements
\usepackage{needspace} % for avoiding page brake right after the section title
\usepackage{iftex} % check if engine is pdflatex, xetex or luatex

% Ensure that generate pdf is machine readable/ATS parsable:
\ifPDFTeX
    \input{glyphtounicode}
    \pdfgentounicode=1
    \usepackage[T1]{fontenc}
    \usepackage[utf8]{inputenc}
    \usepackage{lmodern}
\fi

\usepackage{charter}

% Some settings:
\raggedright
\AtBeginEnvironment{adjustwidth}{\partopsep0pt} % remove space before adjustwidth environment
\pagestyle{empty} % no header or footer
\setcounter{secnumdepth}{0} % no section numbering
\setlength{\parindent}{0pt} % no indentation
\setlength{\topskip}{0pt} % no top skip
\setlength{\columnsep}{0.15cm} % set column seperation
\pagenumbering{gobble} % no page numbering

\titleformat{\section}{\needspace{4\baselineskip}\scshape\raggedright\large}{}{0pt}{}[\vspace{1pt}\titlerule]

\titlespacing{\section}{
    % left space:
    -1pt
}{
    % top space:
    0.15 cm
}{
    % bottom space:
    0.1 cm
} % section title spacing

\renewcommand\labelitemi{$\vcenter{\hbox{\small$\bullet$}}$} % custom bullet points
\newenvironment{highlights}{
    \begin{itemize}[
        topsep=0.10 cm,
        parsep=0.10 cm,
        partopsep=0pt,
        itemsep=0pt,
        leftmargin=0 cm + 10pt
    ]
}{
    \end{itemize}
} % new environment for highlights


\newenvironment{highlightsforbulletentries}{
    \begin{itemize}[
        topsep=0.10 cm,
        parsep=0.10 cm,
        partopsep=0pt,
        itemsep=0pt,
        leftmargin=10pt
    ]
}{
    \end{itemize}
} % new environment for highlights for bullet entries

\newenvironment{onecolentry}{
    \begin{adjustwidth}{
        0 cm + 0.00001 cm
    }{
        0 cm + 0.00001 cm
    }
}{
    \end{adjustwidth}
} % new environment for one column entries

\newenvironment{twocolentry}[2][]{
    \onecolentry
    \def\secondColumn{#2}
    \setcolumnwidth{\fill, 4.5 cm}
    \begin{paracol}{2}
}{
    \switchcolumn \raggedleft \secondColumn
    \end{paracol}
    \endonecolentry
} % new environment for two column entries

\newenvironment{threecolentry}[3][]{
    \onecolentry
    \def\thirdColumn{#3}
    \setcolumnwidth{, \fill, 4.5 cm}
    \begin{paracol}{3}
    {\raggedright #2} \switchcolumn
}{
    \switchcolumn \raggedleft \thirdColumn
    \end{paracol}
    \endonecolentry
} % new environment for three column entries

\newenvironment{header}{
    \setlength{\topsep}{0pt}\par\kern\topsep\centering\linespread{1.5}
}{
    \par\kern\topsep
} % new environment for the header

\newcommand{\placelastupdatedtext}{% \placetextbox{<horizontal pos>}{<vertical pos>}{<stuff>}
  \AddToShipoutPictureFG*{% Add <stuff> to current page foreground
    \put(
        \LenToUnit{\paperwidth-2 cm-0 cm+0.05cm},
        \LenToUnit{\paperheight-1.0 cm}
    ){\vtop{{\null}\makebox[0pt][c]{
        \small\color{gray}\textit{Last updated in July 2024}\hspace{\widthof{Last updated in July 2024}}
    }}}%
  }%
}%

% save the original href command in a new command:
\let\hrefWithoutArrow\href

% new command for external links:


\begin{document}
    \newcommand{\AND}{\unskip
        \cleaders\copy\ANDbox\hskip\wd\ANDbox
        \ignorespaces
    }
    \newsavebox\ANDbox
    \sbox\ANDbox{$|$}

    \begin{header}
        \fontsize{25 pt}{25 pt}\selectfont Tianyue Zhou

        \vspace{5 pt}

        \normalsize
        \kern 5.0 pt%
        \mbox{tianyuez@mit.edu, zhoutianyue1@gmail.com}%
        \kern 5.0 pt%
        \AND%
        \kern 5.0 pt%
        \mbox{+86 13913830510, (617)4603149}%
        \kern 5.0 pt%
    \end{header}

    \vspace{5 pt - 0.3 cm}

    \section{Education}
            \begin{twocolentry}{
            Sept 2025 - Present
        }
            \textbf{Massachusetts Institute of Technology (MIT)}\end{twocolentry}
        \begin{twocolentry}{
            \textit{Cambridge, MA}
        }
            \textit{Ph.D. in Civil and Environmental Engineering; Advisor: \href{http://www.wucathy.com/blog/}{\textcolor{blue}{Prof. Cathy Wu}}; GPA:\textbf{5.0/5.0} }\end{twocolentry}
        \vspace{0.10 cm}
        \begin{twocolentry}{Sept 2021 - June 2025}
            \textbf{ShanghaiTech University}
        \end{twocolentry}
        \begin{twocolentry}{
            \textit{Shanghai, China}
        }
            \textit{B.Eng. in Computer Science; GPA: \textbf{3.75/4.0}}
        \end{twocolentry}
        \vspace{0.10 cm}
        
        % \begin{onecolentry}
        %     \begin{highlights}
        %         \item \textbf{Coursework:} Deep Learning (A), Introduction to Machine Learning (A), Numerical Optimization (A), Probability and Statistics (A+), Mathematical Analysis I \& II (A \& A+), Linear Algebra (A), Computational Science and Engineering (A+), Data Structures and Algorithms (A).
        %     \end{highlights}
        % \end{onecolentry}
        \vspace{0.10 cm}
        \begin{twocolentry}{
            Sept 2023 - Dec 2023
        }
            \textbf{Massachusetts Institute of Technology (MIT)}\end{twocolentry}
        \begin{twocolentry}{
            \textit{Cambridge, MA}
        }
            \textit{Undergraduate Exchange Student, Computer Science; GPA: \textbf{4.5/5.0}}\end{twocolentry}
        \vspace{0.10 cm}
        % \begin{onecolentry}
        %     \begin{highlights}
        %         \item \textbf{Coursework:} Reinforcement Learning: Foundations and Methods (A).
        %     \end{highlights}
        % \end{onecolentry}
    \section{Research Interests}
        \begin{highlights}
            \item \textbf{Machine Learning}: Reinforcement Learning, Multi-task Learning, Transfer Learning.
        \end{highlights}

    \section{Publications}
    \begin{highlights}
        \item \textbf{Expert with Clustering: Hierarchical Online Preference Learning
Framework [\href{https://proceedings.mlr.press/v242/zhou24a/zhou24a.pdf}{\textcolor{magenta}{paper}}]}\\
        \textit{\underline{\textbf{Tianyue Zhou}}, Jung-Hoon Cho, Babak Rahimi Ardabili, Hamed Tabkhi, Cathy Wu}\\
        Learning for Dynamics and Control Conference, 707–718
        % Accepted by 6th Annual Learning for Dynamics \& Control Conference (L4DC) 2024
        
        \item \textbf{The Nah Bandit: Modeling User Non-compliance in
 Recommendation Systems [\href{https://ieeexplore.ieee.org/document/11130914}{\textcolor{magenta}{paper}}] [\href{https://tianyuejo.github.io/The-Nah-Bandit/}{\textcolor{magenta}{website}}] [\href{https://github.com/TianyueJo/The-Nah-Bandit}{\textcolor{magenta}{code}}]}\\
        \textit{\underline{\textbf{Tianyue Zhou}}, Jung-Hoon Cho, Cathy Wu}\\
        IEEE Transactions on Control of Network Systems (Volume: 12, Issue: 4, Dec. 2025), 2919-2931

        \item \textbf{Structure Detection for Contextual Reinforcement
Learning [\href{https://ojs.aaai.org/index.php/AAAI/article/view/40137}{\textcolor{magenta}{paper}}] [\href{https://mit-wu-lab.github.io/SD-MBTL/}{\textcolor{magenta}{website}}] [\href{https://github.com/mit-wu-lab/SD-MBTL}{\textcolor{magenta}{code}}]}\\
        \textit{\underline{\textbf{Tianyue Zhou}}, Jung-Hoon Cho, Cathy Wu}\\
        Proceedings of the AAAI Conference on Artificial Intelligence, 40(34), 29009–29016.

    % \item \textbf{Learning-based Incentive Design for Eco-Driving Guidance [\href{https://trc-30.epfl.ch/wp-content/uploads/2024/09/TRC-30_paper_292.pdf}{\textcolor{magenta}{paper}}]}\\
    %     \textit{Jung-Hoon Cho, M. Umar B. Niazi, Siqi Du, \underline{\textbf{Tianyue Zhou}}, Roy Dong, Cathy Wu}\\
    %     30th Anniversary of Transportation Research: Part C (TRC-30)
    \end{highlights}
    
    \section{Research Experience}
    \begin{twocolentry}{\textit{Cambridge, MA}}
        \textbf{Massachusetts Institute of Technology}
    \end{twocolentry}
    \begin{twocolentry}{}
         \textit{Undergraduate Research Assistant, Supervised by \href{http://www.wucathy.com/blog/}{\textcolor{blue}{Prof. Cathy Wu}}, MIT LIDS}
    \end{twocolentry}
    
    \begin{highlights}
        \item \textbf{Online Preference Learning for Drivers} \hfill Sept 2023 - Jan 2024
        \begin{highlightsforbulletentries}
            \vspace{-0.1cm}
            \item \textbf{Objective}: Implement online learning algorithm in a travel route recommendation problem, aiming to rapidly capture user preferences and recommend personalized and eco-friendly travel routes.
            \item Developed a novel contextual bandit framework, Expert with Clustering (EWC), leveraging hierarchical user information for rapid and accurate learning of user preferences, achieving a \textbf{27.57\%} performance improvement over LinUCB baseline in experiments.
            \item Proposed a distance metric, Loss-guided Distance, tailored for the online preference learning problem, which enhances the representativeness of centroids in clustering.
            \item Established the regret bound of EWC, indicating superior
theoretical performance compared to LinUCB.
        \end{highlightsforbulletentries}
        \vspace{-0.1cm}
        
        \item \textbf{Modeling User Non-compliance in Recommendation Systems} \hfill Jan 2024 - June 2025
        \begin{highlightsforbulletentries}
        \vspace{-0.1cm}
        \item \textbf{Objective}: Tackle a critical challenge often overlooked by traditional frameworks in physical-world recommendation problems: users can easily dismiss recommendations that do not appeal to them and revert to their baseline behavior.
        \item Introduced a novel bandit problem--Nah Bandit--for modeling user non-compliance in recommendation systems. This framework offers the potential to accelerate preference learning compared to traditional frameworks.
        \item Proposed a user non-compliance model to parameterize the anchoring effect in Nah Bandit, reducing the bias when learning user preference.
        \item Combined the user non-compliance model with EWC, further efficiently utilizing non-compliance feedback and hierarchical information. Experimental results highlight that EWC achieves \textbf{at least 15.36\%} performance improvement in multiple applications compared to both supervised learning and traditional bandits.
        \end{highlightsforbulletentries}
        \vspace{-0.1cm}
        \item \textbf{Structure Detection for Contextual Reinforcement Learning} \hfill Aug 2024 - Aug 2025
        \begin{highlightsforbulletentries}
        \vspace{-0.1cm}
        \item \textbf{Objective}: Extend model-based transfer learning approach to address multi-dimensional contextual reinforcement learning (cRL) problems.
        \item Proposed SD-MBTL, a generic framework that detects underlying generalization structures and adapts task selection strategies accordingly.
        \item Developed M/GP-MBTL, an initialization of SD-MBTL that combines a clustering based approach and a Gaussian Process-based approach, tailored to different structural patterns in CMDPs.
        \item Validated on multi-dimensional real-world benchmarks, including robotics, eco-driving, and agriculture, showing significant gains in generalization and sample efficiency.
        \end{highlightsforbulletentries}
        \vspace{-0.1cm}
    \end{highlights}
    
    \section{Honors and Awards}
    \begin{highlights}
\item ShanghaiTech \textbf{Outstanding Student} (2021-2022), Top $\mathbf{3\%\sim7\%}$ \hfill \textit{Dec 2022}
\item ShanghaiTech \textbf{Outstanding Student Union Officer} (2021-2022) \hfill \textit{Dec 2022}
\item \textbf{Scholarship} for Undergraduate 3+1 Overseas Exchange Program, $\sim$\$9,000 \hfill \textit{June 2024}
    \end{highlights}
    \section{Miscellaneous}
        \begin{highlights}
            \item \textbf{Programming Languages:} Python, C/C++, MATLAB
            \item \textbf{Tools:} PyTorch, Git, LaTeX, Markdown
        \end{highlights}
        
\end{document}